\begin{trapbox}
\begin{enumerate}
\item \textbf{长轴顶点和短轴顶点的最值写反.}
切线到中心的距离满足 $b\le d\le a$.
长轴顶点 $(\pm a,0)$ 对应 $d=a$ (最大),
短轴顶点 $(0,\pm b)$ 对应 $d=b$ (最小).
原版例 2 的题干 ``短轴顶点处最大'' 必须纠正.

\item \textbf{把 $d$ 错看成 $OP$.}
$d$ 是\textbf{从中心到切线的垂距}, 并不是中心到切点的线段长.
一般情况下切线不垂直于 $OP$;
只有轴端点等特殊位置 (或圆的情形) 才能直接把二者看成同一个长度.

\item \textbf{把焦半径 $r_i$ 与焦点到切线的垂距 $h_i$ 混为一谈.}
$r_i=PF_i$, $h_i=\operatorname{dist}(F_i,l)$.
正确关系是 $h_i=r_i\sin\theta$,
以及 $h_1h_2=b^2$; 不能写成 $r_1r_2=b^2$ 对任意点都成立.

\item \textbf{把几何平均数错换成算术平均数.}
椭圆定义给出 $(r_1+r_2)/2=a$, 但并不意味着
$\sqrt{r_1r_2}=a$.
由均值不等式, $\sqrt{r_1r_2}\le a$;
两者仅在 $r_1=r_2$ (短轴端点) 时相等.

\item \textbf{用 $r_1+r_2=2a$ 推最值时漏掉下界.}
仅由均值不等式得 $r_1r_2\le a^2$, 只能推出 $d\ge b$.
欲推出 $d\le a$, 还需用
$|r_1-r_2|\le 2c$ (三角形不等式), 或 $|x_0|\le a$,
得到 $r_1r_2\ge a^2-c^2=b^2$.

\item \textbf{混淆 $a,b,c,e$ 的位置与定义.}
本节采用长轴在 $x$ 轴上的写法, $a>b>0$,
$c^2=a^2-b^2$, $e=c/a$.
如果长轴在 $y$ 轴上, 焦半径的坐标表达式需换成含 $y_0$,
但以椭圆中心到切线的距离定义 $d$ 时, 主公式仍是
$\sqrt{r_1r_2}\,d=ab$.

\item \textbf{忽略几何条件和正平方根.}
$d,r_1,r_2,a,b$ 均为正数, 故由
$d^2r_1r_2=a^2b^2$ 应取\textbf{正根}.
题目若声称 $r_1r_2<b^2$ 或 $r_1r_2>a^2$, 则与该椭圆上的点不相容.

\item \textbf{误认为椭圆是圆的拉伸, 所以距离直接保留.}
不等比拉伸不会保持欧氏距离或垂直关系.
借助圆进行形象理解可以, 真正证明需要切线方程、反射性质、
投影或三角恒等式, 不可只说 ``由圆类推可得''.

\item \textbf{几何证明中漏掉 ``同侧'' 前提.}
直接使用 $d=(h_1+h_2)/2$ 的理由是两焦点与中心都在切线同侧,
且中心是焦点的中点. 对任意两点、任意直线, 若分处直线两侧,
则必须先讨论有向距离, 不能把两个\textbf{无向}垂距直接取平均.

\item \textbf{把 $ab$ 当成面积.}
$ab$ 是两个半轴长度的乘积, 椭圆面积为 $S=\pi ab$.
若要改用面积表述, 正确写法是
$\sqrt{r_1r_2}\,d=S/\pi$, 而不是 $S$.
\end{enumerate}
\end{trapbox}
